\documentclass[11pt,a4paper]{article}
\usepackage[utf8]{inputenc}
\usepackage[T1]{fontenc}
\usepackage[spanish,es-minimal,es-noshorthands]{babel}
\usepackage{lmodern}
\usepackage[sfdefault,lining]{FiraSans}
\usepackage[a4paper,top=2.3cm,bottom=2.2cm,left=1.8cm,right=1.8cm,headheight=1.3cm,headsep=0.4cm,footskip=1cm]{geometry}
\usepackage{amsmath,amssymb}
\usepackage{icomma}
\usepackage{graphicx}
\usepackage[table]{xcolor}
\usepackage{tikz}
\usepackage{fontawesome5}
\usepackage{ccicons}
\usepackage{fancyhdr}
\usepackage{lastpage}
\usepackage{tasks}
\usepackage{needspace}
\usepackage[most]{tcolorbox}
\usepackage{array}
\usepackage[hidelinks]{hyperref}
\AtBeginDocument{\hypersetup{pdftitle={Soluciones de la ficha de la unidad 2: Divisibilidad. 1º ESO},pdfsubject={Matemáticas 1º ESO},pdfkeywords={divisibilidad, múltiplos, divisores, criterios de divisibilidad, números primos, factores primos, m.c.m., m.c.d.}}}

\definecolor{navy}{HTML}{1B2748}
\definecolor{accent}{HTML}{E8603A}
\definecolor{teal}{HTML}{1F9E90}
\definecolor{soft}{HTML}{EEF1F8}
\definecolor{muted}{HTML}{5D6685}

\setlength{\parindent}{0pt}
\setlength{\parskip}{3pt}
\settasks{label=\alph*), label-format=\bfseries\color{navy}, label-width=1.3em, item-indent=1.8em, column-sep=0.6em, after-item-skip=2pt, before-skip=3pt, after-skip=4pt}

% ---------- capçalera i peu ----------
\newcommand{\llicencia}{\href{https://creativecommons.org/licenses/by-nc-sa/4.0/deed.es}{\ccbyncsa}\ \ Licencia Creative Commons Atribución-NoComercial-CompartirIgual 4.0 Internacional (CC BY-NC-SA 4.0)}
\pagestyle{fancy}
\fancyhf{}
\fancyhead[L]{\raisebox{-0.35\height}{\includegraphics[height=0.95cm]{logo-gva.png}}}
\fancyhead[C]{\small\color{muted}Soluciones de la ficha · Unidad 2: Divisibilidad · 1\textsuperscript{o} ESO}
\fancyhead[R]{\raisebox{-0.35\height}{\includegraphics[height=0.95cm]{logo-vermellar.png}}}
\fancyfoot[L]{\scriptsize\color{muted}\llicencia}
\fancyfoot[R]{\small\color{muted}\thepage}
\renewcommand{\headrulewidth}{0.4pt}
\renewcommand{\headrule}{\color{muted!40}\hrule height \headrulewidth}
\fancypagestyle{primera}{\fancyhf{}\renewcommand{\headrulewidth}{0pt}
  \fancyfoot[L]{\scriptsize\color{muted}\llicencia}
  \fancyfoot[R]{\small\color{muted}\thepage}}

% ---------- nivells i exercicis ----------
\newcommand{\estrelles}[1]{{\color{accent}\ifcase#1\or\faStar\faStar[regular]\faStar[regular]\or\faStar\faStar\faStar[regular]\or\faStar\faStar\faStar\fi}}
\newcounter{exer}
\newcommand{\ex}[1]{\par\needspace{4\baselineskip}\medskip\stepcounter{exer}\noindent{\color{navy}\bfseries\theexer.}\ \ {\small\estrelles{#1}}\ \ }

\newcommand{\seccio}[2]{\par\needspace{10\baselineskip}\bigskip
  \begin{tcolorbox}[colback=navy,colframe=navy,coltext=white,arc=2mm,boxrule=0pt,left=3mm,right=3mm,top=1.2mm,bottom=1.2mm]
  {\large\bfseries #1\quad #2}
  \end{tcolorbox}\vspace{-2pt}}

\newcolumntype{C}[1]{>{\centering\arraybackslash}p{#1}}

% quadre de recordatori (només a la fitxa d'exercicis)
\newcommand{\recorda}[1]{\par\smallskip\begin{tcolorbox}[colback=accent!7,colframe=accent!55,arc=2mm,boxrule=0.5pt,left=2.5mm,right=2.5mm,top=1mm,bottom=1mm]\small{\color{accent}\faLightbulb[regular]}\ \textbf{Recuerda:}\ #1\end{tcolorbox}}
\newcommand{\res}[1]{\boxed{#1}}

% =====================================================================
% Capçalera FentMates de les fitxes: un color per a cada curs.
% Els colors són els de la pàgina inicial (LEVELS a assets/js/materials.js);
% el logo ve de _eines/logo/pdf/ (python3 _eines/logo/genera.py).
%
% Ús, al final del preàmbul de la unitat:
%   % =====================================================================
% Capçalera FentMates de les fitxes: un color per a cada curs.
% Els colors són els de la pàgina inicial (LEVELS a assets/js/materials.js);
% el logo ve de _eines/logo/pdf/ (python3 _eines/logo/genera.py).
%
% Ús, al final del preàmbul de la unitat:
%   % =====================================================================
% Capçalera FentMates de les fitxes: un color per a cada curs.
% Els colors són els de la pàgina inicial (LEVELS a assets/js/materials.js);
% el logo ve de _eines/logo/pdf/ (python3 _eines/logo/genera.py).
%
% Ús, al final del preàmbul de la unitat:
%   \input{../comu/fentmates.tex}
%   \fmcurs{1eso}
% i, al començament del document:
%   \capcalera{Matemàtiques}{1r ESO}{Fitxa d'exercicis · Unitat 2}{Divisibilitat}
% Opcional abans de \input: \def\fmneed{14} (espai mínim abans de cada secció).
% =====================================================================
\makeatletter
\def\fm@logodir{../../logo/pdf}
\providecommand{\fmneed}{10}
\def\fm@cursnom{}

% clau del curs, color i color del text damunt del color
\newcommand{\fm@curs}[3]{\@namedef{fm@color@#1}{#2}\@namedef{fm@text@#1}{#3}}
\fm@curs{1eso}{E6B84A}{navy}
\fm@curs{2eso}{3B82F6}{white}
\fm@curs{3eso}{FF7A50}{white}
\fm@curs{4eso-a}{A78BFA}{white}
\fm@curs{4eso-b}{6D28D9}{white}
\fm@curs{1batx-mi}{9B1C1C}{white}
\fm@curs{1batx-ccss}{EF6B6B}{white}
\fm@curs{2batx-mii}{1E6B47}{white}
\fm@curs{2batx-ccss}{2BB3A3}{white}

\newcommand{\fmcurs}[1]{%
  \@ifundefined{fm@color@#1}{\PackageError{fentmates}{Curs desconegut: #1}{}}{}%
  \edef\fm@tmp{\noexpand\definecolor{curs}{HTML}{\@nameuse{fm@color@#1}}%
               \noexpand\colorlet{curstext}{\@nameuse{fm@text@#1}}}\fm@tmp
  \colorlet{cursfosc}{curs!72!black}%
  \def\fm@logo{\fm@logodir/vertical-#1.pdf}%
  \def\fm@marca{\fm@logodir/marca-#1.pdf}}

% ---------- capçalera de la primera pàgina ----------
% #1 matèria (petit), #2 curs, #3 tipus de document i unitat, #4 títol
\newcommand{\capcalera}[4]{%
  \gdef\fm@cursnom{#2}%
  \par\noindent
  \begin{tcolorbox}[colback=curs,colframe=curs,boxrule=0pt,arc=2.2mm,width=3.5cm,height=2.15cm,
      valign=center,left=3mm,right=2mm,top=1mm,bottom=1mm,nobeforeafter,box align=center]
    \color{curstext}{\footnotesize #1}\par\vspace{1pt}{\fontsize{19}{21}\selectfont\bfseries #2\par}%
  \end{tcolorbox}\hspace{5mm}%
  \begin{minipage}[c]{\dimexpr\linewidth-3.5cm-5mm-2.9cm\relax}
    {\color{cursfosc}\small\bfseries\MakeUppercase{#3}}\par\vspace{2pt}
    {\color{navy}\fontsize{20}{23}\selectfont\bfseries #4\par}
  \end{minipage}\hfill
  \begin{minipage}[c]{2.6cm}\raggedleft\includegraphics[width=2.6cm]{\fm@logo}\end{minipage}\par
  \vspace{2pt}}

% ---------- barres de secció del color del curs ----------
\renewcommand{\seccio}[2]{\par\needspace{\fmneed\baselineskip}\bigskip
  \begin{tcolorbox}[colback=curs,colframe=curs,coltext=curstext,arc=2mm,boxrule=0pt,left=3mm,right=3mm,top=1.2mm,bottom=1.2mm]
  {\large\bfseries #1\quad #2}
  \end{tcolorbox}\vspace{-2pt}}

% ---------- estrelles de dificultat del color del curs (a Batxillerat no se'n posen) ----------
\@ifundefined{estrelles}{}{%
  \renewcommand{\estrelles}[1]{{\color{curs}\ifcase#1\or\faStar\faStar[regular]\faStar[regular]\or\faStar\faStar\faStar[regular]\or\faStar\faStar\faStar\fi}}}

% ---------- capçalera i peu de les pàgines ----------
% primera pàgina: logos institucionals al peu
\fancypagestyle{primera}{\fancyhf{}\renewcommand{\headrulewidth}{0pt}
  \fancyfoot[L]{\raisebox{-0.32\height}{\includegraphics[height=0.78cm]{logo-gva.png}}\hspace{5mm}%
    \raisebox{-0.32\height}{\includegraphics[height=0.78cm]{logo-vermellar.png}}}
  \fancyfoot[R]{\scriptsize\color{muted}\href{https://creativecommons.org/licenses/by-nc-sa/4.0/deed.ca}{\ccbyncsa}\ \ CC BY-NC-SA 4.0 · fentmates.org\qquad\small\thepage}}
% la resta: el símbol a l'esquerra i el curs a la dreta, amb el títol al centre com abans
\fancyhead[L]{\raisebox{-0.3\height}{\includegraphics[height=0.62cm]{\fm@marca}}}
\fancyhead[R]{\small\bfseries\color{cursfosc}\fm@cursnom}
\renewcommand{\headrule}{\color{curs!45}\hrule height \headrulewidth}
\makeatother

%   \fmcurs{1eso}
% i, al començament del document:
%   \capcalera{Matemàtiques}{1r ESO}{Fitxa d'exercicis · Unitat 2}{Divisibilitat}
% Opcional abans de \input: \def\fmneed{14} (espai mínim abans de cada secció).
% =====================================================================
\makeatletter
\def\fm@logodir{../../logo/pdf}
\providecommand{\fmneed}{10}
\def\fm@cursnom{}

% clau del curs, color i color del text damunt del color
\newcommand{\fm@curs}[3]{\@namedef{fm@color@#1}{#2}\@namedef{fm@text@#1}{#3}}
\fm@curs{1eso}{E6B84A}{navy}
\fm@curs{2eso}{3B82F6}{white}
\fm@curs{3eso}{FF7A50}{white}
\fm@curs{4eso-a}{A78BFA}{white}
\fm@curs{4eso-b}{6D28D9}{white}
\fm@curs{1batx-mi}{9B1C1C}{white}
\fm@curs{1batx-ccss}{EF6B6B}{white}
\fm@curs{2batx-mii}{1E6B47}{white}
\fm@curs{2batx-ccss}{2BB3A3}{white}

\newcommand{\fmcurs}[1]{%
  \@ifundefined{fm@color@#1}{\PackageError{fentmates}{Curs desconegut: #1}{}}{}%
  \edef\fm@tmp{\noexpand\definecolor{curs}{HTML}{\@nameuse{fm@color@#1}}%
               \noexpand\colorlet{curstext}{\@nameuse{fm@text@#1}}}\fm@tmp
  \colorlet{cursfosc}{curs!72!black}%
  \def\fm@logo{\fm@logodir/vertical-#1.pdf}%
  \def\fm@marca{\fm@logodir/marca-#1.pdf}}

% ---------- capçalera de la primera pàgina ----------
% #1 matèria (petit), #2 curs, #3 tipus de document i unitat, #4 títol
\newcommand{\capcalera}[4]{%
  \gdef\fm@cursnom{#2}%
  \par\noindent
  \begin{tcolorbox}[colback=curs,colframe=curs,boxrule=0pt,arc=2.2mm,width=3.5cm,height=2.15cm,
      valign=center,left=3mm,right=2mm,top=1mm,bottom=1mm,nobeforeafter,box align=center]
    \color{curstext}{\footnotesize #1}\par\vspace{1pt}{\fontsize{19}{21}\selectfont\bfseries #2\par}%
  \end{tcolorbox}\hspace{5mm}%
  \begin{minipage}[c]{\dimexpr\linewidth-3.5cm-5mm-2.9cm\relax}
    {\color{cursfosc}\small\bfseries\MakeUppercase{#3}}\par\vspace{2pt}
    {\color{navy}\fontsize{20}{23}\selectfont\bfseries #4\par}
  \end{minipage}\hfill
  \begin{minipage}[c]{2.6cm}\raggedleft\includegraphics[width=2.6cm]{\fm@logo}\end{minipage}\par
  \vspace{2pt}}

% ---------- barres de secció del color del curs ----------
\renewcommand{\seccio}[2]{\par\needspace{\fmneed\baselineskip}\bigskip
  \begin{tcolorbox}[colback=curs,colframe=curs,coltext=curstext,arc=2mm,boxrule=0pt,left=3mm,right=3mm,top=1.2mm,bottom=1.2mm]
  {\large\bfseries #1\quad #2}
  \end{tcolorbox}\vspace{-2pt}}

% ---------- estrelles de dificultat del color del curs (a Batxillerat no se'n posen) ----------
\@ifundefined{estrelles}{}{%
  \renewcommand{\estrelles}[1]{{\color{curs}\ifcase#1\or\faStar\faStar[regular]\faStar[regular]\or\faStar\faStar\faStar[regular]\or\faStar\faStar\faStar\fi}}}

% ---------- capçalera i peu de les pàgines ----------
% primera pàgina: logos institucionals al peu
\fancypagestyle{primera}{\fancyhf{}\renewcommand{\headrulewidth}{0pt}
  \fancyfoot[L]{\raisebox{-0.32\height}{\includegraphics[height=0.78cm]{logo-gva.png}}\hspace{5mm}%
    \raisebox{-0.32\height}{\includegraphics[height=0.78cm]{logo-vermellar.png}}}
  \fancyfoot[R]{\scriptsize\color{muted}\href{https://creativecommons.org/licenses/by-nc-sa/4.0/deed.ca}{\ccbyncsa}\ \ CC BY-NC-SA 4.0 · fentmates.org\qquad\small\thepage}}
% la resta: el símbol a l'esquerra i el curs a la dreta, amb el títol al centre com abans
\fancyhead[L]{\raisebox{-0.3\height}{\includegraphics[height=0.62cm]{\fm@marca}}}
\fancyhead[R]{\small\bfseries\color{cursfosc}\fm@cursnom}
\renewcommand{\headrule}{\color{curs!45}\hrule height \headrulewidth}
\makeatother

%   \fmcurs{1eso}
% i, al començament del document:
%   \capcalera{Matemàtiques}{1r ESO}{Fitxa d'exercicis · Unitat 2}{Divisibilitat}
% Opcional abans de \input: \def\fmneed{14} (espai mínim abans de cada secció).
% =====================================================================
\makeatletter
\def\fm@logodir{../../logo/pdf}
\providecommand{\fmneed}{10}
\def\fm@cursnom{}

% clau del curs, color i color del text damunt del color
\newcommand{\fm@curs}[3]{\@namedef{fm@color@#1}{#2}\@namedef{fm@text@#1}{#3}}
\fm@curs{1eso}{E6B84A}{navy}
\fm@curs{2eso}{3B82F6}{white}
\fm@curs{3eso}{FF7A50}{white}
\fm@curs{4eso-a}{A78BFA}{white}
\fm@curs{4eso-b}{6D28D9}{white}
\fm@curs{1batx-mi}{9B1C1C}{white}
\fm@curs{1batx-ccss}{EF6B6B}{white}
\fm@curs{2batx-mii}{1E6B47}{white}
\fm@curs{2batx-ccss}{2BB3A3}{white}

\newcommand{\fmcurs}[1]{%
  \@ifundefined{fm@color@#1}{\PackageError{fentmates}{Curs desconegut: #1}{}}{}%
  \edef\fm@tmp{\noexpand\definecolor{curs}{HTML}{\@nameuse{fm@color@#1}}%
               \noexpand\colorlet{curstext}{\@nameuse{fm@text@#1}}}\fm@tmp
  \colorlet{cursfosc}{curs!72!black}%
  \def\fm@logo{\fm@logodir/vertical-#1.pdf}%
  \def\fm@marca{\fm@logodir/marca-#1.pdf}}

% ---------- capçalera de la primera pàgina ----------
% #1 matèria (petit), #2 curs, #3 tipus de document i unitat, #4 títol
\newcommand{\capcalera}[4]{%
  \gdef\fm@cursnom{#2}%
  \par\noindent
  \begin{tcolorbox}[colback=curs,colframe=curs,boxrule=0pt,arc=2.2mm,width=3.5cm,height=2.15cm,
      valign=center,left=3mm,right=2mm,top=1mm,bottom=1mm,nobeforeafter,box align=center]
    \color{curstext}{\footnotesize #1}\par\vspace{1pt}{\fontsize{19}{21}\selectfont\bfseries #2\par}%
  \end{tcolorbox}\hspace{5mm}%
  \begin{minipage}[c]{\dimexpr\linewidth-3.5cm-5mm-2.9cm\relax}
    {\color{cursfosc}\small\bfseries\MakeUppercase{#3}}\par\vspace{2pt}
    {\color{navy}\fontsize{20}{23}\selectfont\bfseries #4\par}
  \end{minipage}\hfill
  \begin{minipage}[c]{2.6cm}\raggedleft\includegraphics[width=2.6cm]{\fm@logo}\end{minipage}\par
  \vspace{2pt}}

% ---------- barres de secció del color del curs ----------
\renewcommand{\seccio}[2]{\par\needspace{\fmneed\baselineskip}\bigskip
  \begin{tcolorbox}[colback=curs,colframe=curs,coltext=curstext,arc=2mm,boxrule=0pt,left=3mm,right=3mm,top=1.2mm,bottom=1.2mm]
  {\large\bfseries #1\quad #2}
  \end{tcolorbox}\vspace{-2pt}}

% ---------- estrelles de dificultat del color del curs (a Batxillerat no se'n posen) ----------
\@ifundefined{estrelles}{}{%
  \renewcommand{\estrelles}[1]{{\color{curs}\ifcase#1\or\faStar\faStar[regular]\faStar[regular]\or\faStar\faStar\faStar[regular]\or\faStar\faStar\faStar\fi}}}

% ---------- capçalera i peu de les pàgines ----------
% primera pàgina: logos institucionals al peu
\fancypagestyle{primera}{\fancyhf{}\renewcommand{\headrulewidth}{0pt}
  \fancyfoot[L]{\raisebox{-0.32\height}{\includegraphics[height=0.78cm]{logo-gva.png}}\hspace{5mm}%
    \raisebox{-0.32\height}{\includegraphics[height=0.78cm]{logo-vermellar.png}}}
  \fancyfoot[R]{\scriptsize\color{muted}\href{https://creativecommons.org/licenses/by-nc-sa/4.0/deed.ca}{\ccbyncsa}\ \ CC BY-NC-SA 4.0 · fentmates.org\qquad\small\thepage}}
% la resta: el símbol a l'esquerra i el curs a la dreta, amb el títol al centre com abans
\fancyhead[L]{\raisebox{-0.3\height}{\includegraphics[height=0.62cm]{\fm@marca}}}
\fancyhead[R]{\small\bfseries\color{cursfosc}\fm@cursnom}
\renewcommand{\headrule}{\color{curs!45}\hrule height \headrulewidth}
\makeatother

\fmcurs{1eso}

\begin{document}
\thispagestyle{primera}

\capcalera{Matemáticas}{1\textsuperscript{o} ESO}{Soluciones de la ficha · Unidad 2}{Divisibilidad}
\vspace{8pt}

\begin{tcolorbox}[colback=soft,colframe=soft,arc=2mm,left=3mm,right=3mm,top=1.5mm,bottom=1.5mm]
\small Intenta resolver cada ejercicio antes de mirar su solución. Si no te sale el mismo resultado, repasa los pasos o consulta los apuntes interactivos de la unidad.
\end{tcolorbox}

% =====================================================
\seccio{1}{La relación de divisibilidad}

\ex{1} \begin{tasks}(2)
  \task $84=7\cdot12$: \res{\text{sí}}
  \task $95=8\cdot11+7$: \res{\text{no}}
  \task $132=12\cdot11$: \res{\text{sí}}
  \task $200=15\cdot13+5$: \res{\text{no}}
\end{tasks}

\ex{1} \begin{tasks}(4)
  \task múltiplo
  \task divisor
  \task múltiplo
  \task divisor
\end{tasks}

\ex{1} $72$ es múltiplo de $8$; \ $72$ es múltiplo de $9$; \ $8$ es divisor de $72$; \ $9$ es divisor de $72$.

\ex{2} \begin{tasks}(1)
  \task \res{\text{V}}: $48=6\cdot8$.
  \task \res{\text{V}}: $77=7\cdot11$.
  \task \res{\text{F}}: $120=9\cdot13+3$.
  \task \res{\text{V}}: $0=15\cdot0$.
  \task \res{\text{F}}: $1:3$ no es exacta (es $3$ el que es múltiplo de $1$).
\end{tasks}

\ex{2} $100=9\cdot11+1$. \ Sumar: $\res{8}$ ($108=9\cdot12$). \ Restar: $\res{1}$ ($99=9\cdot11$).

\ex{3} Exacta: $13\cdot17=\res{221}$. \ El resto como mucho es $12$: $221+12=\res{233}$.

\ex{3} De los múltiplos de 7 de dos cifras ($14, 21, \ldots, 98$), acaban en 4: $\res{14}$ y $\res{84}$.


% =====================================================
\seccio{2}{Múltiplos de un número}

\ex{1} \begin{tasks}(1)
  \task $7,\ 14,\ 21,\ 28,\ 35,\ 42,\ 49,\ 56,\ 63,\ 70$
  \task $11,\ 22,\ 33,\ 44,\ 55,\ 66,\ 77,\ 88,\ 99,\ 110$
  \task $15,\ 30,\ 45,\ 60,\ 75,\ 90,\ 105,\ 120,\ 135,\ 150$
\end{tasks}

\ex{1} $\res{18,\ 36,\ 54,\ 60,\ 96}$ \ (los demás no dan división exacta entre 6).

\ex{2} $100=13\cdot7+9$: el primero es $13\cdot8=104$. \ $\res{104,\ 117,\ 130,\ 143,\ 156,\ 169}$.

\ex{2} $200:9\to q=22$: \ hay $\res{22}$. \ Hasta $99$ hay $11$, así que entre 100 y 200: $22-11=\res{11}$ (del $108$ al $198$).

\ex{2} $\res{18,\ 36,\ 54,\ 72,\ 90}$. El más pequeño es $18$ (su m.c.m.).

\ex{3} $300=17\cdot17+11$: los múltiplos más próximos son $289$ (a 11) y $306$ (a 6). \ El más próximo es $\res{306}$.

\ex{3} $999=12\cdot83+3$: \ $12\cdot83=\res{996}$ \ y \ $12\cdot84=\res{1008}$.


% =====================================================
\seccio{3}{Divisores de un número}

\ex{1} \begin{tasks}(2)
  \task $\{1,\ 3,\ 5,\ 15\}$
  \task $\{1,\ 3,\ 7,\ 21\}$
  \task $\{1,\ 2,\ 4,\ 7,\ 14,\ 28\}$
  \task $\{1,\ 2,\ 4,\ 8,\ 16,\ 32\}$
\end{tasks}

\ex{1} $\res{4,\ 12,\ 15,\ 30}$ \ ($60=4\cdot15=12\cdot5=15\cdot4=30\cdot2$).

\ex{2} \begin{tasks}(1)
  \task $\{1,2,3,4,6,8,12,16,24,48\}$
  \task $\{1,2,3,4,6,8,9,12,18,24,36,72\}$
  \task $\{1,2,3,5,6,9,10,15,18,30,45,90\}$
  \task $\{1,2,4,5,10,20,25,50,100\}$
\end{tasks}

\ex{2} $\res{1,\ 2,\ 3,\ 4,\ 6,\ 12}$. El más grande es $12$ (su m.c.d.).

\ex{2} Tantas como divisores tiene 36: $1\times36,\ 2\times18,\ 3\times12,\ 4\times9,\ 6\times6,\ 9\times4,\ 12\times3,\ 18\times2,\ 36\times1$: \ $\res{9\text{ maneras}}$.

\ex{3} El $\res{48}$, con 10 divisores: $1, 2, 3, 4, 6, 8, 12, 16, 24, 48$.

\ex{3} Por ejemplo, $\res{6}$: $\{1,2,3,6\}$ \ y \ $\res{16}$: $\{1,2,4,8,16\}$.


% =====================================================
\seccio{4}{Criterios de divisibilidad}

\ex{1} \par\smallskip{\renewcommand{\arraystretch}{1.5}\centering\begin{tabular}{|C{2.4cm}|C{1.5cm}|C{1.5cm}|C{1.5cm}|C{1.5cm}|C{1.5cm}|}\hline\rowcolor{soft}\textbf{Número} & $\boldsymbol{\div2}$ & $\boldsymbol{\div3}$ & $\boldsymbol{\div5}$ & $\boldsymbol{\div9}$ & $\boldsymbol{\div11}$\\\hline $3080$ & \checkmark & $\times$ & \checkmark & $\times$ & \checkmark\\\hline $7425$ & $\times$ & \checkmark & \checkmark & \checkmark & \checkmark\\\hline $1188$ & \checkmark & \checkmark & $\times$ & \checkmark & \checkmark\\\hline $2543$ & $\times$ & $\times$ & $\times$ & $\times$ & $\times$\\\hline $6735$ & $\times$ & \checkmark & \checkmark & $\times$ & $\times$\\\hline \end{tabular}\par}\par El $2543$ no es divisible por ninguno: es impar, las cifras suman 14 y la diferencia para el 11 es $8-6=2$.

\ex{1} Por ejemplo, $\res{102,\ 111,\ 123}$: las cifras suman 3, 3 y 6 (múltiplos de 3, pero no de 9).

\ex{2} $5+4=9$. \ Por 3: $\square=0,3,6,9$ \ ($\res{504,\ 534,\ 564,\ 594}$). \ Por 9: $\square=0$ o $9$ \ ($\res{504,\ 594}$).

\ex{2} Por 3: $11+\square$ múltiplo de 3, \ $\square=1, 4, 7$. Por 2 tiene que ser par: $\square=4$, \ $\res{384}$.

\ex{2} \begin{tasks}(2)
  \task $(9+8)-(4+2)=11$: \res{\text{sí}}
  \task $(7+9)-(5+3)=8$: \res{\text{no}}
  \task $(2+1+5)-(8+0)=0$: \res{\text{sí}}
  \task $(5+5+1)-(7+3)=1$: \res{\text{no}}
\end{tasks}

\ex{3} Tiene que acabar en 0 (por 2 y por 5) y las cifras tienen que sumar un múltiplo de 3: $\res{990}$ ($9+9+0=18$).

\ex{3} Por 5: $b=0$ o $5$. \ Si $b=0$: $9+a$ múltiplo de 9, $a=0$ o $9$. \ Si $b=5$: $14+a$, $a=4$. \ \ $\res{7020,\ 7920,\ 7425}$.


% =====================================================
\seccio{5}{Números primos y compuestos}

\ex{1} \begin{tasks}(3)
  \task $\{1,11\}$: primo
  \task $\{1,3,5,15\}$: compuesto
  \task $\{1,19\}$: primo
  \task $\{1,3,9,27\}$: compuesto
  \task $\{1,31\}$: primo
  \task $\{1,7,49\}$: compuesto
\end{tasks}

\ex{1} $\res{61,\ 67,\ 71,\ 73,\ 79,\ 83,\ 89,\ 97}$.

\ex{2} \begin{tasks}(2)
  \task $87=3\cdot29$: \res{\text{compuesto}}
  \task $107:11\to q=9<11$: \res{\text{primo}}
  \task $133=7\cdot19$: \res{\text{compuesto}}
  \task $149:13\to q=11<13$: \res{\text{primo}}
\end{tasks}

\ex{2} \begin{tasks}(2)
  \task $3+13=5+11$
  \task $5+19=7+17=11+13$
  \task $5+31=7+29=13+23=17+19$
  \task $3+47=7+43=13+37=19+31$
\end{tasks}\par (Basta con una manera.)

\ex{2} Si acaba en 5, es divisible por 5: tiene al menos los divisores 1, 5 y él mismo. Si es par, es divisible por 2: tiene al menos los divisores 1, 2 y él mismo.

\ex{3} $2027$ no es divisible por ningún primo hasta el $43$, y $47\cdot47=2209>2027$: es \res{\text{primo}}. \ $2028$ es par: \res{\text{compuesto}} ($2028=2^{2}\cdot3\cdot13^{2}$).

\ex{3} Por ejemplo, $5+7+11=\res{23}$ \ o \ $7+11+13=\res{31}$.


% =====================================================
\seccio{6}{Descomposición en factores primos}

\ex{1} \begin{tasks}(4)
  \task $2^{2}\cdot3^{2}$
  \task $2\cdot5^{2}$
  \task $2^{3}\cdot3^{2}$
  \task $3^{4}$
\end{tasks}

\ex{1} \begin{tasks}(4)
  \task $2^{3}\cdot3\cdot5$
  \task $2\cdot3^{2}\cdot7$
  \task $2^{2}\cdot7^{2}$
  \task $3^{2}\cdot5^{2}$
\end{tasks}

\ex{2} \begin{tasks}(4)
  \task $2^{2}\cdot3^{3}\cdot5$
  \task $2^{4}\cdot7^{2}$
  \task $2^{2}\cdot3^{2}\cdot5\cdot7$
  \task $2\cdot3\cdot5\cdot7\cdot11$
\end{tasks}

\ex{2} \begin{tasks}(4)
  \task $48$
  \task $126$
  \task $500$
  \task $165$
\end{tasks}

\ex{2} \begin{tasks}(3)
  \task $2^{3}\cdot3\cdot5^{2}$
  \task $2^{2}\cdot3^{2}\cdot7$
  \task $2\cdot3^{2}\cdot5$
\end{tasks}

\ex{3} $90=2\cdot3^{2}\cdot5$: falta un 2 y un 5 para que todos los exponentes sean pares. Por $\res{10}$: $900=2^{2}\cdot3^{2}\cdot5^{2}=30^{2}$.

\ex{3} Es el $60$. \ Divisores: $\res{1,\ 2,\ 3,\ 4,\ 5,\ 6,\ 10,\ 12,\ 15,\ 20,\ 30,\ 60}$.


% =====================================================
\seccio{7}{Mínimo común múltiplo}

\ex{1} \begin{tasks}(4)
  \task $12$
  \task $18$
  \task $40$
  \task $15$
\end{tasks}

\ex{1} \begin{tasks}(2)
  \task $2^{2}\cdot3\cdot5=\res{60}$
  \task $2\cdot3^{2}\cdot5=\res{90}$
  \task $2^{3}\cdot3^{2}=\res{72}$
  \task $2\cdot3\cdot7=\res{42}$
\end{tasks}

\ex{2} \begin{tasks}(2)
  \task $2^{3}\cdot3\cdot5=\res{120}$
  \task $2^{2}\cdot3^{3}=\res{108}$
  \task $3^{2}\cdot5^{2}=\res{225}$
  \task $2^{2}\cdot3\cdot7=\res{84}$
\end{tasks}

\ex{2} \begin{tasks}(3)
  \task $2^{3}\cdot3=\res{24}$
  \task $2^{2}\cdot3\cdot5=\res{60}$
  \task $2^{3}\cdot3^{3}=\res{216}$
\end{tasks}

\ex{2} \begin{tasks}(2)
  \task $\res{90}$: no tienen factores comunes
  \task $\res{35}$: 35 es múltiplo de 7
  \task $\res{143}$: los dos son primos
  \task $\res{48}$: 48 es múltiplo de 16
\end{tasks}

\ex{3} $2^{3}\cdot3^{3}\cdot5\cdot7=8\cdot27\cdot35=\res{7560}$.

\ex{3} $60=2^{2}\cdot3\cdot5$ y $12=2^{2}\cdot3$: el otro tiene que dividir a 60 y tener el factor 5. \ $\res{5,\ 10,\ 15,\ 20,\ 30,\ 60}$.


% =====================================================
\seccio{8}{Problemas de m.c.m.}

\ex{1} $\text{m.c.m.}(9,12)=2^{2}\cdot3^{2}=\res{36\text{ s}}$.

\ex{1} $\text{m.c.m.}(3,4)=\res{12\text{ días}}$.

\ex{2} $\text{m.c.m.}(6,8,10)=2^{3}\cdot3\cdot5=\res{120\text{ días}}$.

\ex{2} $\text{m.c.m.}(12,15,18)=2^{2}\cdot3^{2}\cdot5=\res{180\text{ lápices}}$.

\ex{2} $\text{m.c.m.}(20,25)=2^{2}\cdot5^{2}=100$ min $=1$ h $40$ min: \ a las $\res{9{:}10}$.

\ex{3} $\text{m.c.m.}(6,8,10)=120$: como mínimo $120+3=\res{123}$. \ Entre 200 y 300: $240+3=\res{243}$.

\ex{3} $\text{m.c.m.}(15,24)=2^{3}\cdot3\cdot5=120$ dientes: \ la primera da $120:15=\res{8\text{ vueltas}}$ y la segunda, $120:24=\res{5\text{ vueltas}}$.


% =====================================================
\seccio{9}{Máximo común divisor}

\ex{1} \begin{tasks}(4)
  \task $4$
  \task $5$
  \task $9$
  \task $1$
\end{tasks}

\ex{1} \begin{tasks}(2)
  \task $2^{2}\cdot3=\res{12}$
  \task $3\cdot5=\res{15}$
  \task $2\cdot7=\res{14}$
  \task $2\cdot3\cdot5=\res{30}$
\end{tasks}

\ex{2} \begin{tasks}(2)
  \task $2\cdot3\cdot7=\res{42}$
  \task $2^{2}\cdot3\cdot5=\res{60}$
  \task $2^{4}\cdot3=\res{48}$
  \task $3\cdot5=\res{15}$
\end{tasks}

\ex{2} \begin{tasks}(3)
  \task $2\cdot3=\res{6}$
  \task $2^{3}=\res{8}$
  \task $3\cdot5=\res{15}$
\end{tasks}

\ex{2} \begin{tasks}(2)
  \task \res{\text{sí}} ($\text{m.c.d.}=1$)
  \task no ($\text{m.c.d.}=5$)
  \task \res{\text{sí}} ($\text{m.c.d.}=1$)
  \task no ($\text{m.c.d.}=11$)
\end{tasks}

\ex{3} $36=2^{2}\cdot3^{2}$, \ $60=2^{2}\cdot3\cdot5$: \ $\text{m.c.d.}=\res{12}$, \ $\text{m.c.m.}=\res{180}$. \ $12\cdot180=2160=36\cdot60$ \checkmark

\ex{3} $\text{m.c.d.}\cdot\text{m.c.m.}=a\cdot b$: \ $8\cdot96=768$; \ $768:24=\res{32}$.


% =====================================================
\seccio{10}{Problemas de m.c.d.}

\ex{1} $\text{m.c.d.}(30,45)=15$: \ $\res{15\text{ lotes}}$, cada uno con $2$ manzanas y $3$ naranjas.

\ex{1} $\text{m.c.d.}(36,48)=\res{12\text{ m}}$; \ $3+4=\res{7\text{ trozos}}$.

\ex{2} $\text{m.c.d.}(360,280)=2^{3}\cdot5=\res{40\text{ cm}}$; \ $9\cdot7=\res{63\text{ baldosas}}$.

\ex{2} $\text{m.c.d.}(120,84,60)=\res{12\text{ libros}}$; \ $10+7+5=\res{22\text{ paquetes}}$.

\ex{2} \begin{tasks}(1)
  \task m.c.m.: $\text{m.c.m.}(15,25)=\res{75\text{ min}}$.
  \task m.c.d.: $\text{m.c.d.}(48,64)=\res{16\text{ estuches}}$, con 3 bolígrafos y 4 lápices cada uno.
\end{tasks}

\ex{3} Se han repartido $72$ cromos y $96$ pegatinas: \ $\text{m.c.d.}(72,96)=\res{24\text{ alumnos}}$ (más grande que 9). Cada uno recibe $\res{3}$ cromos y $\res{4}$ pegatinas.

\ex{3} $\text{m.c.d.}(72,48)=\res{24\text{ m}}$. \ Perímetro: $2\cdot(72+48)=240$ m; \ $240:24=\res{10\text{ árboles}}$.


\end{document}
